\section{Despliegue y confianza}

\subsection{Revisión del despliegue y reversión}

Después de cada release hay que hacer pre-mortem, post-mortem y verificación de rollback:
\begin{itemize}
    \item \textbf{Pre-mortem}: antes del rollout, simular el worst-case prompt drift y confirmar la cobertura de los guardrails.
    \item \textbf{Post-mortem}: recopilar el rejection log y el reward score drift, y anotarlos en el incident log.
    \item \textbf{Verificación de rollback}: diseñar un rollback playbook y ensayar primero en staging el restore checkpoint y el recover guardrail.
\end{itemize}

\begin{importantbox}{El valor de anticipar incidentes}
Hacer un simulacro de estimación de riesgos antes de que el gerente de producto apruebe el rollout permite incluir los high-risk prompts en una watchlist y evitar improvisar con prisa cuando ocurra un incidente.
\end{importantbox}

\begin{table}[H]
\centering
\begin{tabular}{p{0.25\textwidth} p{0.65\textwidth}}
\toprule
Etapa de revisión & Contenido \\
\midrule
Pre-mortem & drift risk list, revisión de la evidence matrix, guardrail traceability \\
Post-mortem & reward score drift graph, human override log, root cause analysis \\
Rollback & reference checkpoint + runbook, alert escalation path \\
\bottomrule
\end{tabular}
\caption{Entregables de la revisión del despliegue y la reversión}
\end{table}

\subsection{Gobernanza con múltiples roles y confianza}

Un alignment pipeline no es solo research: el equipo de gobernanza, el área legal y SRE también deben participar:
\begin{itemize}
    \item El equipo de investigación y desarrollo se encarga del reward model y del policy update.
    \item El equipo de operación y monitoreo se encarga de los drift alerts y del dashboard.
    \item El área legal y de cumplimiento se encarga del guardrail audit y de la privacy review.
    \item El PM y el decision board deciden en última instancia el rollout o el rollback con base en el evidence ranking.
\end{itemize}

\begin{knowledgebox}{Propuesta de gobernanza con múltiples roles}
Declarar la evidence matrix como un multi-stakeholder checkpoint; por ejemplo, en cada policy update, el compliance team debe revisar de nuevo los long-form prompts.
\end{knowledgebox}

\begin{warningbox}{Cadena de amplificación del riesgo}
Si falta un governance checkpoint, puede ocurrir que, después del fine-tune del reward model, se eleve cierto tipo de prompt sin que compliance lo detecte (catch), de modo que las high-risk replies lleguen directamente a production.
\end{warningbox}

\subsection{Resumen de la sección}

El despliegue y la confianza requieren pre-mortem/post-mortem+rollback y una evidence governance con múltiples roles; solo así se puede mantener la confianza en el alignment dentro de un RLHF pipeline en iteración continua.
